\section*{Chapter \ref{ch:g7:functions} --- Functions}

\begin{solution}{exo:g7:functions:1}
\code{Hello!}, then \code{between}, then \code{Bye!}: the body of \code{shout}
runs at each call, in the order of the \omterm{def:g4:programs-and-operating-system:program}{program}.
\end{solution}

\begin{solution}{exo:g7:functions:2}
$f(2) = 5$ and $f(3) = 7$, so it prints 12.
\end{solution}

\begin{solution}{exo:g7:functions:3}
\code{doubling 5}, printed during the call, then \code{11}: the call returned
10, kept in \code{result}, and 1 was added.
\end{solution}

\begin{solution}{exo:g7:functions:4}
\omcode{code/grade-7/01-functions/exercises/perimeter.py}{1}{6}{The perimeter of
a rectangle.}\label{lst:g7:functions:sol-perimeter}

It prints 14 and 21.0.
\end{solution}

\begin{solution}{exo:g7:functions:5}
\code{average} prints the average but returns nothing, so \code{average(8, 12)}
is \code{None}, and \code{None * 2} is refused. Replace the \code{print} by
\code{return (a + b) / 2}; the \omterm{def:g4:programs-and-operating-system:program}{program} then prints 20.0.
\end{solution}

\begin{solution}{exo:g7:functions:6}
No: \code{area(4, 3)} gives 12, like \code{area(3, 4)}, because $4 \times 3 = 3
\times 4$. For \code{introduce}, the two \omterm{def:g7:functions:parameter}{parameters} play different parts in the
sentence, so swapping the \omterm{def:g7:functions:parameter}{arguments} gives a wrong sentence. The order of the
\omterm{def:g7:functions:parameter}{arguments} always decides which \omterm{def:g7:functions:parameter}{parameter} gets which \omterm{def:g6:variables-and-input:value}{value}; it only happens not
to change the result of \code{area}.
\end{solution}

\begin{solution}{exo:g7:functions:7}
\omcode{code/grade-7/01-functions/exercises/is_even.py}{1}{5}{A \omterm{def:g7:functions:function}{function} that
returns a \omterm{def:g6:if-else:comparison}{Boolean value}.}\label{lst:g7:functions:sol-is-even}

The \omterm{def:g6:if-else:comparison}{comparison} is itself \code{True} or \code{False}: no \code{if} is needed. It
prints \code{True False True}.
\end{solution}

\begin{solution}{exo:g7:functions:8}
\omcode{code/grade-7/01-functions/exercises/minutes_to_hours.py}{1}{6}{Minutes
into hours.}\label{lst:g7:functions:sol-minutes}

Because it returns its \omterm{def:g6:variables-and-input:value}{value}, the results can be added: $0.75 + 0.25 = 1.0$.
\end{solution}

\begin{solution}{exo:g7:functions:9}
4 times, once for each turn of the \omterm{def:g4:algorithms-repeat-and-decide:loop}{loop}: it prints \code{computing 1} to
\code{computing 4}, then $1 + 4 + 9 + 16 = 30$.
\end{solution}

\begin{solution}{exo:g7:functions:10}
Each branch returns its price, so the \omterm{def:g7:functions:function}{function} ends as soon as a case matches.

\omcode{code/grade-7/01-functions/exercises/price.py}{1}{11}{The ticket price as
a \omterm{def:g7:functions:function}{function}.}\label{lst:g7:functions:sol-price}

It prints 5, 9, 9 and 6 for 11, 12, 64 and 65.
\end{solution}

\begin{solution}{exo:g7:functions:11}
Three calls put the largest number in each place: \code{largest3(9, 2, 4)},
\code{largest3(2, 9, 4)} and \code{largest3(2, 4, 9)}, all expected to give 9.
They reach both branches of each call to \code{largest2}. Then the equal cases:
\code{largest3(5, 5, 1)}, \code{largest3(1, 5, 5)}, \code{largest3(5, 1, 5)}
and \code{largest3(5, 5, 5)}, all expected to give 5. That makes at least
seven calls.
\end{solution}

\begin{solution}{exo:g7:functions:12}
\omcode{code/grade-7/01-functions/exercises/digit_sum.py}{1}{14}{The digit sum,
checked on five examples.}\label{lst:g7:functions:sol-digit-sum}

The examples cover an ordinary number, zero (the \omterm{def:g4:algorithms-repeat-and-decide:loop}{loop} does not run at all), a
single digit, zeros inside a number, and a long number.
\end{solution}

\begin{solution}{pb:g7:functions:1}
\textbf{1.} From the right, the digits 7, 5, 4, 3, 2, 1, 4, 3, 2, 1, 0, 9, 5
are multiplied by 1, 3, 1, 3, \dots: $7 + 15 + 4 + 9 + 2 + 3 + 4 + 9 + 2 + 3 + 0
+ 27 + 5 = 90$, a multiple of 10.

\textbf{2.} With a check digit 0, the weighted sum would be 89. The check digit
has weight 1, so it must add 1: the check digit is 1, and the number is
\code{4006381333931}.

\textbf{3.} The second digit from the right is 8 instead of 5. Its weight is 3,
so the sum grows by $3 \times 3 = 9$: 99, not a multiple of 10. Refused.

\textbf{4.} With weight 1, the check digit is simply the amount missing to
reach the next multiple of 10. With weight 3, we would have to find a digit
whose triple fills the gap.

\textbf{5.} When \code{weight} is 1, \code{4 - weight} is 3; when it is 3, it
is 1. So the weights alternate.

\textbf{6.} Multiplying by 10 adds a 0 on the right: a place for the check
digit, with weight 1, and every other digit moves to the weight it will have in
the full number. If the sum ends in 0, \code{10 - 0} gives 10, and the second
\texttt{\% 10} turns it into the digit 0.

\textbf{7.} \code{check_digit(400638133393)} returns 1; \code{is_valid} returns
\code{True} for \code{5901234123457} and \code{False} for \code{5901234123487}.

\textbf{8.} Work out the answers by hand first, then print them beside the
calls:
\begin{itemize}
  \item \code{check_digit(590123412345)}: expected 7 (question 1);
  \item \code{check_digit(400638133393)}: expected 1 (question 2);
  \item \code{check_digit(100000000003)}: expected 0; the sum is already a
    multiple of 10 ($1 + 3 \times 3 = 10$), which tests the second \texttt{\% 10};
  \item \code{check_digit(100000000000)}: expected 9 (the sum is 1).
\end{itemize}

\textbf{9.} 13 places with 9 other digits each: 117 numbers. All 117 are
refused:

\omoutput{code/grade-7/01-functions/outputs/one_digit.txt}{Every single wrong
digit is caught.}\label{out:g7:functions:one-digit}

\textbf{10.} Changing a digit by $d$ (between $-9$ and 9, not 0) changes the
sum by $3 \times d$. For that to be a multiple of 10, $3 \times d$ would have to
end in 0; but $3 \times 1, \dots, 3 \times 9$ are 3, 6, 9, 12, 15, 18, 21, 24,
27: none ends in 0. The same holds for negative $d$. With weight 1, the change
$d$ itself is between $-9$ and 9 and not 0, so it is not a multiple of 10
either.

\textbf{11.} Before: $a \times 1 + b \times 3$; after: $b \times 1 + a \times
3$. The change is $2 \times a - 2 \times b = 2 \times (a - b)$.

\textbf{12.} 80 of the 90 swaps are caught:

\omoutput{code/grade-7/01-functions/outputs/swaps.txt}{The 10 missed swaps.}\label{out:g7:functions:swaps}

$2 \times (a - b)$ is a multiple of 10 only when $a - b$ is 5 or $-5$: the
missed swaps are exactly the pairs of digits that differ by 5.

\textbf{13.} \textbf{Result.} Every single wrong digit is caught. Of the 90
swaps of two different neighbouring digits, 80 are caught, and the 10 that are
missed are exactly those where the two digits differ by 5.
% ledger: g7:ean13-errors
\end{solution}
